\documentclass[10pt,twocolumn,a4paper]{article}
\usepackage[margin=1.8cm,columnsep=0.6cm]{geometry}
\usepackage[T1]{fontenc}
\usepackage{lmodern}
\usepackage{amsmath,amssymb}
\usepackage{graphicx}
\usepackage{booktabs}
\usepackage{multirow}
\usepackage{xcolor}
\usepackage{algorithm}
\usepackage{algpseudocode}
\usepackage{caption}
\usepackage[hidelinks]{hyperref}
\usepackage{titlesec}
\titleformat{\section}{\normalfont\large\bfseries}{\thesection}{0.6em}{}
\titleformat{\subsection}{\normalfont\normalsize\bfseries}{\thesubsection}{0.6em}{}
\captionsetup{font=small,labelfont=bf}
\setlength{\tabcolsep}{4pt}
\newcommand{\M}{\texttt{[M]}}

\title{\textbf{Can a Tabular Diffusion Model Generate Drivable Trips?}\\[2pt]
\large TabDiff on Road-Traversal Transactions, with a Graph-Aware Sampler}
\author{Gradient Gladiators: Abhay Sharma and Naveen Kumar\\
\small IIIT Hyderabad \quad DFMGAI Course Project}
\date{}

\begin{document}
\maketitle

\begin{abstract}
TabDiff (ICLR 2025) is a continuous-time diffusion model for mixed-type tables. We first reproduce it on
the Adult benchmark, matching the published Trend error (1.46\% against 1.49\%). We then ask whether it can
generate road trips. Following Rathan, Reddy and Mondal (DAPD 2022), we turn taxi trajectories from T-Drive
(Beijing) into market-basket transactions whose items are the road segments travelled, giving a
300-column all-categorical table of 42,705 trips. TabDiff keeps which roads are busy (Spearman 0.81) but
only 61\% of its trips are drivable paths. We propose a graph-aware re-masking sampler that sends
segments violating road connectivity back to the mask during sampling. It raises drivable trips to 87\%,
and a matched random control shows the gain comes from the road graph, not from re-masking as such. The
cost is shorter trips and weaker busy-road agreement. We compare against Markov, independent and
DiffTraj baselines and discuss what this says about tabular diffusion for structured data.
\end{abstract}

%======================================================================
\section{Problem Statement and Motivation}
\label{sec:a}
Trajectory data from vehicles drives traffic planning, routing and map services, but it is sensitive: a
trip reveals where a person lives and works. Synthetic trips that keep the statistics of the real ones,
without copying them, would let such data be shared.

Rathan, Reddy and Mondal~\cite{reddy2022} showed that road-network traversals can be mined with
frequent-pattern methods once each trip is written as a \emph{transaction}: a set of road segments. Popular
paths then follow from frequent itemsets. A table of transactions is exactly the kind of data a tabular
generative model learns. TabDiff~\cite{tabdiff} is the current state of the art for tabular synthesis.

\textbf{Question.} Can TabDiff generate road-traversal transactions that (i) are drivable paths on the road
graph, and (ii) preserve the busy roads, frequent patterns and popular paths of the real data? And
(iii) does a sampler that knows the road graph improve on TabDiff's default sampler?

%======================================================================
\section{Gaps Addressed}
\label{sec:b}
\begin{itemize}\setlength\itemsep{1pt}
\item \textbf{Tabular models ignore structure between columns.} TabDiff, like other tabular generators,
  treats columns as exchangeable features. Road segments are not: a valid trip is a connected path, a
  hard constraint that no column-wise score (Shape, Trend) measures.
\item \textbf{No evaluation of tabular diffusion on transaction data.} TabDiff was evaluated on seven UCI
  tables with at most a few dozen columns, mostly dense. Road transactions have 300 sparse binary columns
  (about 4 segments per trip are ``on'').
\item \textbf{Sampling that can undo mistakes.} In masked diffusion a decoded value is final. TabDiff's
  stochastic sampler can re-noise, but it is switched off for all-categorical tables, and it is blind to
  domain constraints. We add a re-masking rule driven by the road graph.
\end{itemize}

%======================================================================
\section{Problem Formulation}
\label{sec:c}
\textbf{Transactions.} Let $G=(V,E)$ be a directed road graph and $E_k\subset E$ the $k$ most used
segments in a study window. A trip is map-matched to a sequence of segments; its transaction is the
binary vector $x\in\{0,1\}^{k}$ with $x_j=1$ if segment $j$ is travelled. A transaction is
\emph{drivable} if its on-segments form one simple directed path in $G$ (no gaps, no branches).

\textbf{Generation.} Given a training table $\mathcal{D}=\{x^{(n)}\}$, learn $p_\theta(x)$ and sample a
synthetic table $\tilde{\mathcal{D}}$ such that drivability is high and the following match the real test
split: segment-use counts (busy roads), frequent itemsets (patterns), and the ranking of alternative
paths between the same origin and destination (popular paths, Reddy et al.'s score $\omega$).

\textbf{TabDiff's forward process.} Each categorical column $j$ with $C_j$ classes is extended with an
absorbing mask class \M. In continuous time $t\in[0,1]$, with a learnable per-column schedule
$\alpha_j(t)$ decreasing from 1 to 0,
\begin{equation}
q(x_{t,j}\mid x_{0,j}) = \alpha_j(t)\,\delta_{x_{0,j}} + \big(1-\alpha_j(t)\big)\,\delta_{\M},
\end{equation}
where $\delta_c$ puts all mass on class $c$. Numerical columns receive Gaussian noise with schedule $\sigma_j(t)$; our table has none.

%======================================================================
\section{Algorithm}
\label{sec:d}
\textbf{TabDiff.} A transformer $\mu_\theta(x_t,t)$ predicts the clean value of every masked column.
Sampling starts from the all-mask row at $t=1$ and runs $T=50$ steps to $t=0$. At each step from $t$ to
$s<t$ a masked column is unmasked with probability $\frac{\alpha(s)-\alpha(t)}{1-\alpha(t)}$, drawing its
class from $\mu_\theta$; unmasked columns are copied.

\textbf{Graph-aware re-masking (ours).} Before each denoising step we inspect the decided on-segments of
every row and flag those that break drivability: segments in a branch (two on-segments leave or enter
the same node), boundary segments of a row that splits into several pieces, and off-segments that would
bridge two pieces. Flagged columns are re-masked with extra probability $\lambda$ (the strength), so the
model can decide them again with more context. Training is unchanged; a trained checkpoint is reused.

\begin{algorithm}[t]
\caption{Graph-aware re-masking sampler}\label{alg:gs}
\small
\begin{algorithmic}[1]
\Require trained $\mu_\theta$, road incidence $U,V$, strength $\lambda$
\State $x \gets$ all \M
\For{$i=T,\dots,1$}
  \State $F \gets \textsc{Flag}(x;U,V)$ \Comment{branch, boundary, bridge}
  \State $m \gets m_{\text{TabDiff}} + (1-m_{\text{TabDiff}})\,\lambda\,F$
  \State re-mask each decided column with probability $m$
  \State unmask with $\mu_\theta(x,t_i)$ (TabDiff step)
\EndFor
\State \Return $x$
\end{algorithmic}
\end{algorithm}

\textbf{Control.} To separate the effect of the graph from that of re-masking, a \emph{random} arm
re-masks, in every row and step, as many decided columns as the graph rule flags there, chosen
uniformly at random.

%======================================================================
\section{Theory}
\label{sec:e}
\textbf{Training objective.} For masked diffusion the negative ELBO reduces to a weighted cross-entropy on
masked positions~\cite{tabdiff}:
\begin{equation}
\mathcal{L}_{\text{cat}} = \mathbb{E}_{t,x_0,x_t}\Big[\sum_{j:\,x_{t,j}=\M}
\tfrac{\alpha_j'(t)}{1-\alpha_j(t)}\log \langle \mu_{\theta,j}(x_t,t),\,x_{0,j}\rangle\Big],
\end{equation}
where $\alpha_j'(t)<0$. TabDiff learns the schedule parameters of each column jointly with $\theta$.

\textbf{Absorbing state.} \M\ can be entered but not left in the forward process, so in the reverse process
a decoded value is never revisited. An early wrong guess (a segment that cannot be joined to the rest of
the trip) stays. Re-masking breaks this: with re-masking probability $m$, the reverse chain becomes a
predictor--corrector scheme whose stationary target is still $p_\theta$, but which revisits values.

\textbf{Why the graph-aware rule shortens trips.} A trip has about 4 on-segments out of 300 columns, so
for any column the model's marginal is strongly ``off''. A flagged on-segment that is re-masked is decided
again from a distribution dominated by ``off''. Re-masking therefore removes more segments than it
repairs: validity rises because invalid segments disappear, not because missing ones are filled. This
predicts shorter rows and a loss of busy-road agreement, which is what we observe (Section~\ref{sec:i}).

%======================================================================
\section{Experiments and Reproduction}
\label{sec:f}
We reproduced TabDiff on Adult (32,561 rows; 6 numerical, 9 categorical columns) with the authors' code
and settings, training 8,000 epochs. Table~\ref{tab:repro} compares the final checkpoint with the paper.

\begin{table}[h]
\centering
\caption{Reproduction on Adult: error rates (\%, lower is better).}
\label{tab:repro}
\begin{tabular}{lcc}
\toprule
 & Shape error & Trend error \\
\midrule
TabDiff (paper) & $0.63\pm0.05$ & $1.49\pm0.16$ \\
TabDiff (ours) & 0.76 & 1.46 \\
\midrule
TabSyn (paper) & 0.81 & 1.93 \\
\bottomrule
\end{tabular}
\end{table}

Our Trend error matches the paper; Shape is 0.13 points worse, still better than the strongest baseline
reported in the paper (TabSyn). The reproduction confirms the code runs as published before we move to a
new kind of data.

%======================================================================
\begin{table*}[t]
\centering
\caption{T-Drive test split (8,761 trips). Means over 3 seeds with 95\% $t$-intervals where the method is stochastic. Higher is better except for gaps and copies.}
\label{tab:main}
\begin{tabular}{lcccccccc}
\toprule
Method & Drivable & Gaps & Seg./row & Spearman & Top-25 & Itemset Jaccard & Path $\tau$ & Copies \\
\midrule
Real test (reference) & 1.000 & 0.000 & 4.26 & 0.87 & 0.80 & 0.60 & 0.92 & 0.99 \\
Independent & $0.058\pm.004$ & 0.92 & 4.28 & $0.91\pm.03$ & 0.83 & 0.00 & 0.50 & 0.00 \\
Markov & $1.000$ & 0.00 & 4.31 & $0.93\pm.03$ & 0.85 & $0.51\pm.29$ & 0.74 & 0.93 \\
\midrule
TabDiff, default & $0.610\pm.010$ & 0.37 & 3.86 & $0.81\pm.03$ & 0.48 & $0.10\pm.05$ & 0.85 & 0.47 \\
TabDiff, random 0.1 (control) & $0.626\pm.011$ & 0.35 & 3.85 & $0.81\pm.02$ & 0.48 & $0.11\pm.05$ & 0.85 & 0.49 \\
TabDiff, graph-aware 0.1 & $0.870\pm.010$ & 0.12 & 2.57 & $0.52\pm.06$ & 0.23 & $0.03\pm.01$ & 0.50 & 0.62 \\
\midrule
DiffTraj, fine-tuned & $0.838\pm.001$ & 0.08 & 3.86 & $0.62\pm.06$ & 0.48 & $0.34\pm.11$ & 0.50 & 0.63 \\
DiffTraj, from scratch & $0.708\pm.018$ & 0.14 & 3.70 & $0.58\pm.03$ & 0.31 & $0.02\pm.04$ & 0.69 & 0.49 \\
\bottomrule
\end{tabular}
\end{table*}

\section{Extensions Done in the Project}
\label{sec:g}
\textbf{Data pipeline.} T-Drive~\cite{tdrive}: 10,357 Beijing taxis, 2--8 February 2008, about 15 million
points sampled every $\sim$3 minutes. Road graph from OpenStreetMap. Steps:
(1) clean points (duplicates, speed spikes) and cut trips at stops;
(2) map-match a random sample of 1,000 taxis with an HMM matcher (leuvenmapmatching), 98.6\% of points
matched; (3) pick a 4\,km window in central Beijing and keep its $k=300$ most used directed segments;
(4) write each trip as a 300-column binary transaction. Result: 42,705 trips with 4.3 segments on
average, split \emph{by taxi} into train / validation / test (29,378 / 4,566 / 8,761), so no taxi appears in
two splits.

\textbf{Training.} TabDiff with the authors' architecture, batch 1,024 (11\,GB GPU), 2,000 epochs on one RTX
2080 Ti (17.9 GPU hours). The checkpoint (epoch 750) and the sampler strength ($\lambda=0.1$) were chosen on
validation only.

\textbf{Baselines.}
\emph{Independent}: each segment drawn from its training frequency.
\emph{Markov}: a first-order chain over segments on the road graph, fitted on training trips (drivable by
construction).
\emph{DiffTraj}~\cite{difftraj}, a diffusion model for GPS trajectories: fine-tuned from the authors'
Chengdu checkpoint (50 epochs) and trained from scratch (200 epochs), on the same trips; its output points
are map-matched back to segments.

\textbf{Metrics.} \emph{Validity}: share of drivable rows, share with gaps, segments per row.
\emph{Busy roads}: Spearman correlation of segment-use counts, top-25 overlap. \emph{Patterns}: Jaccard of
frequent multi-segment itemsets (FP-growth, support threshold minF = third quartile of single-segment
support), contiguity of maximal patterns. \emph{Popular paths}: Kendall $\tau$ between real and synthetic
rankings of alternative paths for 26 origin--destination queries. \emph{Privacy}: share of rows that are
exact copies of a training row. TabDiff and DiffTraj use 3 sampling seeds; we report means with 95\%
$t$-intervals.

%======================================================================
\section{New Results and Algorithms}
\label{sec:h}
\textbf{Graph-aware sampler.} Algorithm~\ref{alg:gs} is our algorithmic contribution. It needs only the
node--segment incidence matrices $U,V\in\{0,1\}^{k\times|V|}$: the counts of decided on-segments leaving and
entering each node are $xU$ and $xV$, so all three flags are computed with two matrix products per step,
on the GPU, for the whole batch.

\textbf{Ablation.} Table~\ref{tab:abl} isolates the effect of the graph. Random re-masking touches about
three times as many segments per step as the graph rule, yet leaves TabDiff almost unchanged
(+1.6 points). The graph-aware rule adds 26 points of drivable rows. So the gain comes from
\emph{where} the model re-masks, not from re-masking as such.

\begin{table}[h]
\centering
\caption{Ablation at epoch 750, strength 0.1, seeds 10--12 (test, mean $\pm$ s.d.).}
\label{tab:abl}
\small
\begin{tabular}{lccc}
\toprule
Arm & Re-masked & Drivable (\%) & Seg./row \\
\midrule
A. None & 0 & $61.0\pm0.3$ & 3.86 \\
B. Random & 0.49\% & $62.6\pm0.3$ & 3.85 \\
C. Graph & 0.16\% & $87.0\pm0.3$ & 2.57 \\
\bottomrule
\end{tabular}
\end{table}

%======================================================================
\section{New Experimental Results}
\label{sec:i}


\begin{figure}[t]
\centering
\includegraphics[width=\columnwidth]{figures/tradeoff.pdf}
\caption{Validity against busy-road agreement on the test split (error bars: 95\% intervals over seeds).
The graph-aware sampler moves TabDiff right (more drivable) and down (less faithful).}
\label{fig:tradeoff}
\end{figure}

Table~\ref{tab:main} and Figure~\ref{fig:tradeoff} give the main results.

\textbf{TabDiff keeps marginals, not structure.} With the default sampler, segment frequencies are close
to real (Spearman 0.81) and popular-path rankings are good ($\tau=0.85$), but only 61\% of rows are drivable
and 37\% have gaps. Frequent multi-segment patterns are mostly lost (Jaccard 0.10 against 0.60 for real
data): the model learns each column well but not which columns go together.

\textbf{The graph-aware sampler trades fidelity for validity.} It raises drivable rows to 87\% and makes
every maximal pattern contiguous, but rows shrink from 3.9 to 2.6 segments and busy-road agreement drops
to 0.52, as predicted in Section~\ref{sec:e}.

\textbf{Baselines.} The Markov chain, valid by construction, beats every generative model on fidelity but
copies 93\% of training rows. DiffTraj fine-tuned is more drivable than default TabDiff (84\%) and keeps more
patterns (0.34), but loses busy-road agreement (0.62). Trained from scratch it is clearly worse:
pre-training matters for it. No generative model reaches both high validity and high fidelity.

\begin{figure}[t]
\centering
\includegraphics[width=\columnwidth]{figures/itemsets.pdf}
\caption{Frequent multi-segment itemsets recovered (Jaccard with the real training set's itemsets).}
\label{fig:items}
\end{figure}

\textbf{Privacy.} Exact copies are high even for real test trips (99\%): in a 4\,km window many trips use the
same few segments. Against that reference TabDiff copies less (47\%) than any other method that keeps
structure.

%======================================================================
\section{Conclusion and Limitations}
TabDiff reproduces on Adult and transfers to road transactions in a limited way: it learns which roads are
used, but treats a trip as a bag of columns and breaks connectivity in four rows out of ten. A graph-aware
re-masking sampler, which costs two matrix products per step and no retraining, raises validity by 26
points, and a matched control shows the road graph is what helps. Its weakness is structural: with sparse
rows, re-deciding a segment usually turns it off. Constraining the \emph{choice} of segments during
decoding, rather than re-masking after the fact, is the natural next step.

\textbf{Limitations.} One city, one 4\,km window and one training run; popular-path queries cluster around
few destinations, and with minF = 423 Reddy et al.'s $\omega$ is 0 for every candidate path, so we also
report $\omega$ without the minF rule. T-Drive is sparsely sampled, so map-matching fills in segments between
GPS points.

\begin{thebibliography}{9}\small
\bibitem{tabdiff} J.~Shi, M.~Xu, H.~Hua, H.~Zhang, S.~Ermon, J.~Leskovec. TabDiff: a mixed-type diffusion model
for tabular data generation. \emph{ICLR}, 2025.
\bibitem{reddy2022} P.~R.~Rathan, P.~K.~Reddy, A.~Mondal. A framework for discovering popular paths using
transactional modeling and pattern mining. \emph{Distributed and Parallel Databases}, 40:109--133, 2022.
\bibitem{tdrive} J.~Yuan, Y.~Zheng, C.~Zhang, W.~Xie, X.~Xie, G.~Sun, Y.~Huang. T-Drive: driving directions
based on taxi trajectories. \emph{ACM SIGSPATIAL GIS}, 2010.
\bibitem{difftraj} Y.~Zhu, Y.~Ye, S.~Zhang, X.~Zhao, J.~Yu. DiffTraj: generating GPS trajectory with diffusion
probabilistic model. \emph{NeurIPS}, 2023.
\end{thebibliography}

\end{document}
